% Level 3, Week 11 -- Project Studio 2
\lesson{11}{Project Studio 2}{Test it, playtest it, polish it, explain it.}
\lessoninfo{Concepts: pytest for your own classes, playtesting, bug triage, README docs, accessibility and usability}{Tools: Thonny, pytest, Terminal + git}{CSTA: 3A-AP-19, 3A-AP-21, 3A-AP-23, 2-AP-18}

\begin{objectives}
  \item The student can write at least five pytest tests for their own classes, including an edge case. (L3.G8)
  \item The student can run a playtest, collect feedback and sort bugs by how much they matter (triage). (L3.G5, L3.G8)
  \item The student can write a README that explains what the project does, how to run it, how it works and who to credit. (L3.G5)
  \item The student can review their project for accessibility and usability and fix at least one problem. (L3.G9)
\end{objectives}

\begin{materials}
  \item Print: workbook chapter 11, homework 11, and the showcase talk card page from workbook chapter 12 (Exercise 12.1, needed for homework 11); 8 playtest feedback cards per student (cut from workbook Exercise 11.2, or plain index cards).
  \item On each Mac: the student's project folder open in Thonny; Terminal open in the same folder (\texttt{cd \textasciitilde/Documents/coding-projects/<name>}).
  \item Check pytest works in Thonny: \menu{Tools} $\rightarrow$ \menu{Open system shell...}, type \texttt{python -m pytest --version}. (If it is missing: \menu{Tools} $\rightarrow$ \menu{Manage packages...}, search \texttt{pytest}, \menu{Install}. No account needed.)
  \item Prep (10 min): read this plan and the Project teacher guide's Milestones 3 and 4.
\end{materials}

\begin{timeplan}
  \timeslot{0--6}{Warm-up}{Commit count: everyone runs \texttt{git log --oneline} and reads their best commit message aloud.}
  \timeslot{6--14}{Concept}{What makes a good test; README sections; accessibility and usability checklist.}
  \timeslot{14--22}{Guided demo}{You write three tests for a tiny Player class and run pytest; show one failing test and fix it.}
  \timeslot{22--52}{Hands-on build}{M3: five tests (10 min). Playtest swap + bug triage (8 min). M4: README, accessibility fixes, polish (12 min). Commit after each.}
  \timeslot{52--60}{Wrap-up}{Showcase talk card, exit ticket, homework (finish M4, practice the talk).}
\end{timeplan}

\section{Warm-up (0--6 min)}
Everyone runs \texttt{git log --oneline} and counts commits (target by the showcase: 8+), then reads their best message aloud; vote for the clearest. Only 1--2 commits is fine: today has four natural commit points.

\section{Concept (6--14 min)}
\begin{teacherscript}
\enquote{Today your project grows up. Three jobs.

\textbf{Tests} are tiny programs that check your program. Test the \emph{logic}, not the drawing: \emph{does the player stop at walls? does eating a star add 10 points?} Test normal cases and \textbf{edge cases} -- the corners, the empty list, the last life.

\textbf{Playtesting}: someone else uses your project while you watch silently and take notes. It feels hard not to help -- but every time they get stuck, you learn something.

\textbf{Documentation}: a README tells a stranger what it is, how to run it, how it works and who you thank. Future-you is a stranger too!}
\end{teacherscript}

\begin{analogy}[Car crash tests]
Car makers crash test cars with dummies before anyone drives them. Tests are your crash dummies: you can break things on purpose, safely, as often as you like. Playtesters are the first real passengers.
\end{analogy}

\begin{teachernote}[The CS idea, for you]
Tests are easiest when logic lives in a file with no turtle code. If logic is tangled with drawing, help move one function out and test that. Five small tests beat one giant test; each name says what it checks.
\end{teachernote}

\section{Guided demo (14--22 min)}
\begin{workedexample}[Three tests for a tiny Player]
Use the example in workbook chapter 11 (\enquote{Tests for a tiny Player}): a \texttt{Player} class whose \texttt{move} method stops at walls and scores 10 for a star, plus three tests. Type both files into a new folder (or let a student read them to you), then in Thonny's system shell run \texttt{python -m pytest -q}: three green dots.
\begin{enumerate}[itemsep=1pt]
  \item Read each test name aloud: it says what it checks. Point out the two \textbf{edge cases}: walking into a wall, eating the same star twice.
  \item Break the code on purpose: delete the line \texttt{grid[new\_row][new\_col] = "."} and run again. The third test fails, and its name says exactly what broke.
  \item Put the line back: green again. Commit: \texttt{git commit -m "Add tests for Player"}.
\end{enumerate}
\end{workedexample}

\section{Hands-on build (22--52 min)}
\begin{pairroles}[Builder and tester]
Each student works on their own project; partners sit together. During test writing, the partner is the \textbf{tester}: suggests one edge case (\enquote{what if the list is empty?}). During the playtest, roles become \textbf{player} and \textbf{silent note-taker}. With 1 student: you are the playtester. With 3: play in a circle (A tests B, B tests C, C tests A). With 4: two pairs, then swap pairs for a second playtest if time allows.
\end{pairroles}

\textbf{M3, part 1 -- tests (10 min).} At least five tests in a \texttt{test\_<project>.py} file; at least one edge case. Run pytest; commit when green.

\begin{unplugged}[Playtest with feedback cards (8 min)]
Swap seats. The player uses the partner's project for 3 minutes, thinking aloud. The owner says nothing and writes one feedback card per problem or delight: \emph{What happened? What did they expect?} Then swap.
\end{unplugged}

\begin{bughunt}[Playtest bug triage]
Each owner sorts their cards into three piles: \textbf{Must fix} (crashes, cannot win, cannot start), \textbf{Should fix} (confusing, unfair, hard to read), \textbf{Nice to have} (new ideas). Fix one \enquote{must fix} now. For each crash, first write a test that shows the bug, then fix it -- the test makes sure it never comes back. Commit with a message like \texttt{Fix crash when the player walks off the edge}.
\end{bughunt}

\textbf{M4 -- README and polish (12 min).} Write \texttt{README.md} using workbook Exercise 11.3 (what it is, how to run, controls, how it works -- classes, data structure, algorithm --, tests, credits). Then do the accessibility check below and fix one thing.

\begin{freebuild}[Polish time]
Use any spare minutes to add one \enquote{juice} detail: a title screen, a win message, a high-score table, a nicer palette. Commit each one.
\end{freebuild}

\section{Wrap-up (52--60 min)}
\begin{safety}[Accessibility and usability review]
Every student checks their project against this list (also in the workbook): \textbf{colors} from the Okabe--Ito palette and never color alone (add shapes, letters or text); \textbf{text} large enough to read from a chair; \textbf{controls} written on screen; no fast flashing; \textbf{usability}: a new player knows what to do within 10 seconds. \textbf{Credits}: every borrowed image or sound is listed with its license (or there are none).
\end{safety}

\begin{formativecheck}[Milestones 3 and 4 check and exit ticket]
\textbf{Look for (M3):} five or more passing tests including an edge case; feedback cards sorted. \textbf{(M4):} a README with all sections; at least one accessibility fix; 6+ commits.
\textbf{Question:} \enquote{Name one edge case you tested and why it matters.}
\textbf{Answer:} any real edge, e.g. \enquote{walking into the corner wall -- players will try it, and it used to crash}.
\end{formativecheck}

\begin{reflection}
What did your playtester do that surprised you? Which test are you proudest of? What will you say first in your 2-minute showcase talk?
\end{reflection}

Hand out homework 11 (finish M4, fill in the showcase talk card, practice once).

\section{Notes for the teacher}
\begin{commonmistakes}
  \item \textbf{Tests import turtle code} and open windows: move logic into a file without \texttt{import turtle}, or guard drawing with \py|if __name__ == "__main__":|.
  \item \textbf{Test file not found}: pytest only collects files named \texttt{test\_*.py} and functions named \texttt{test\_*}.
  \item \textbf{Tests that cannot fail} (\py|assert True|, or no \texttt{assert}). Ask: \enquote{Break the code -- does this test go red?}
  \item \textbf{Owner helps during the playtest}: remind them the silence is the point.
\end{commonmistakes}

\begin{differentiation}
\supportpath Give a test template: copy the guided-demo tests and change only the class name and numbers. Three tests today, two more as homework. README: write one sentence for each item of the workbook checklist (Exercise 11.3).
\extensionpath Test the algorithm directly (e.g. BFS finds the shortest path in a tiny 3-by-3 maze); add a \enquote{How it works} diagram to the README.
\end{differentiation}

\begin{stretchchallenge}[Test the save file]
Write a test that saves a high score to a temporary file and loads it back (pytest's \texttt{tmp\_path} fixture gives a safe folder: \py|def test_save(tmp_path):|).
\end{stretchchallenge}

\begin{teachernote}[If you are short on time]
Keep the tests and the playtest; move the README to homework (it is the M4 homework anyway). Shorten the demo to the first two tests.
\end{teachernote}
